% ===========================================================================
\section{What the judges use}
\label{sec:mechanism}

The classifier shows that the frontier texts carry enough signal to name
their authors at \EngineAAcc{}. The frontier judges reach at most
\GPTOverallLineup{} overall. We test four accounts of what they use instead.

\subsection{Rationales}
\label{sec:rationales}
Every frontier judgment came with a one-sentence reason. We coded all 1,900
with a transparent rule-based scheme of eleven cue types (length,
punctuation and formatting, structure, sentence construction, word choice,
tone, hedging, rhetorical flourish, appeals to a model's typical style,
references to the judge itself, and content). \RationaleValidation{}

\emph{No judge ever refers to itself.} Not one of the 1,900 reasons says
anything like ``this is how I write'', including the 66 in which Claude
correctly named itself. Judges instead appeal to what a model typically
writes: 46\% of lineup reasons and 95\% of single-text reasons invoke a
model's characteristic style (``classic Claude hedging'').

\emph{The cited cues are the same whether a self-naming is right or wrong.}
When Claude names itself on its own text in the single-text format it cites
hedging 31 times in 33; when it names itself on a text it did not write, 12
times in 13. GPT's self-namings cite word choice in 33 of 36 hits and 64 of
79 false alarms. A judge that recognised its own writing would not describe
its misses in the same terms as its hits; a judge applying a stereotype of
itself would.

\emph{The cues are mostly true of the text, but not diagnostic.} When a
lineup reason calls a response long, that response is the longest of the
five in 47 of 71 cases (chance 20\%), and phrases judges quote appear
verbatim in 238 of 245 cases. The exception is ``sophisticated vocabulary'',
which matches the response with the longest words no more often than chance
(3 of 35). Judges describe the texts accurately and then map the description
onto a model by reputation. Citing hedging predicts a correct attribution
(odds ratio 1.96, $p<10^{-4}$, lineup), but that shrinks once the true
author is controlled for (1.54, $p=0.017$), because hedging is simply common
in Claude's text. As \citet{marioriyad2025silent} warn, none of this shows
what drove a decision; it shows what judges say drove it.

\subsection{Generative match}
\label{sec:genmatch}
\citet{wataoka2024perplexity} propose that a judge favours, and perhaps
claims, text it would itself have produced. With closed models we can
approximate ``what $J$ would write'' by $J$'s own answer to the same prompt,
and ask whether a text's similarity to it predicts $J$ naming itself on text
it did not write. It does not: across three similarity measures (stylometric
distance, word and character $n$-gram cosine, function-word profile), two
formats and five judges, 1 of 36 clustered models reaches $p<0.05$. GPT, the
judge with the most such false alarms (79 in single-text), shows odds ratio
1.52 [0.85, 2.71] per standard deviation of stylometric similarity. Nor do
judges name the model whose answer to the prompt is most similar to the
text: in 630 wrong lineup attributions the named model is the most similar
one 31.9\% of the time against 30.1\% for shuffled guesses ($p=0.12$),
although the same similarity measures identify the true author 74--79\% of
the time. The signal a generative-match strategy would need is present; the
judges' errors do not follow it. Section~\ref{sec:likelihood} tests the
account directly, with each open-weight judge's own probabilities.

\subsection{Likelihood}
\label{sec:likelihood}
\citet{wataoka2024perplexity} tie self-preference to familiarity: a judge
favours text that is likely under its own distribution, which its own text
usually is. Open weights let us read that familiarity directly. For every
judge and every text we compute the mean token log-probability of the text
as the judge's reply to its prompt, and subtract the mean of the other three
judges' values for the same text, so that a text predictable to every model
does not count as familiar to one.

\emph{Each judge's probabilities identify its own text.} Picking, for each
prompt, the text a judge finds relatively most likely names its own text on
\OWLikCondAccMin{} of prompts for every judge (chance 25\%), and on \OWLikUncAccMin{}--\OWLikUncAccMax{} even without the
prompt. This is expected, since each text was sampled from its author, but
it means that the information needed to recognise one's own text is present
in every judge's forward pass.

\emph{The judges' answers do not use it.} If familiarity drove self-naming,
a judge would claim the texts of others that it finds relatively likely more
often than those it does not. We test this within each prompt, on texts the
judge did not write, comparing the standardised relative likelihood of the
texts it named as its own with the rest (within-prompt permutation). Across
the \OWFamTests{} judge--format pairs with enough claims to test (lineup,
single text and yes/no; \OWFamUntestable{} pairs claim everything or nearly
nothing), none survives Holm correction; the smallest $p$ is
\OWFamMinP{} (\OWFamMinWho{}; Holm \OWFamMinHolm{}), and there the judge
claims the texts it finds \emph{less} likely. Asked who wrote a text, a judge
answers without consulting the one signal that would let it know.

\subsection{Quality}
\label{sec:quality}
\citet{stamand2026sgtr} show that a judge can appear to recognise itself by
claiming whichever text it rates best. Each open-weight judge rated every
text in both corpora for how well it carries out its task, on a 1--9 scale
read as the expected value over the judge's probabilities for the nine
digits.

\emph{Frontier self-naming is not a quality effect.} With the four
open-weight judges as raters independent of the frontier judges, we ask two
questions of Study~1. Do frontier judges claim the better texts? Among texts
a judge did not write, the ones it named as its own are rated no higher than
the rest for any judge in the lineup (within-prompt permutation, smallest
$p=$~\QFLineupClaimMinP{}, Gemini, whose claims lean towards worse texts);
one text at a time, GPT, the one judge that often claims others' text,
claims the \emph{worse} ones (standardised difference \QFGPTSingleClaimDiff{},
$p=$~\QFGPTSingleClaimP{}). Does self-advantage survive holding quality
fixed? In a GEE of correctness on whether the judge wrote the text and the
text's within-prompt quality, the self term stays significant for Claude
(odds ratio \QFClaudeLineupOR{}, $p=$~\QFClaudeLineupORP{}), GPT
(\QFGPTLineupOR{}, $p=$~\QFGPTLineupORP{}) and Grok (\QFGrokLineupOR{},
$p=$~\QFGrokLineupORP{}). The raters agree only weakly (mean pairwise
Spearman \QFRaterRho{}), since almost every frontier text is rated near the
top of the scale (author means \QFMeanLo{}--\QFMeanHi{}), so this measures
quality coarsely.

\emph{The open-weight judges do not prefer their own text either.} We
measure self-preference as how much more a judge rates its own text than its
peers rate it, minus how much more it rates others' text than its peers do,
so that a generous rater is not counted as self-preferring. It ranges from
\OWQualSelfPrefMin{} to \OWQualSelfPrefMax{} points on the nine-point scale,
and every prompt-bootstrap interval includes zero (largest: Mistral,
\OWMistralQualSelfPref{} \OWMistralQualSelfPrefCI{}). Judges that cannot tell
their text from others' do not favour it.

% ===========================================================================
\subsection{Not a watermark}
\label{sec:watermark}

Could a judge be reading a watermark its vendor planted? Three facts say no.

\emph{A keyed watermark cannot be read in context.} Deployed text watermarks
shift token probabilities by a pseudorandom function of a secret key and the
preceding tokens \citep{kirchenbauer2023watermark, dathathri2024synthid}.
Detection needs the key; the pattern is statistically invisible without it,
by design. Nothing in our protocol gives a judge a detector.

\emph{The dates do not fit.} Our frontier corpus was written between 15 June
and 31 July 2026. Anthropic announced keyed watermarking of Claude text in
August 2026, applied at launch to models released from 2 August and
retrofitted to older models, Opus 4.7 among them, without a stated start date
\citep{anthropic2026watermark, anthropic2026marking}. The judge with the
largest self-advantage wrote its texts before any announced Claude watermark.
Conversely, text from the Gemini app is watermarked with SynthID-Text
\citep{deepmind2026synthid}, and Gemini is one of the judges that does
\emph{not} recognise itself.

\emph{Authorship is identifiable without one.} We generated the open-weight
corpus ourselves from public weights, with no sampling modification, so it
carries no watermark; the stylometric classifier still attributes it at
\OWClfPooled{} (Section~\ref{sec:openweight}). The signal that would let a
judge recognise its own text is stylistic and needs no watermark, and the
open-weight judges, which could not be reading one, do not recognise
themselves either.

